% BHU PYQ -- Manually transcribed & error-corrected
% Paper: BCF-215 : Corporate Laws
% Exam: B.Com. (Hons.) FMM Semester III Examination, 2013-14

\documentclass[11pt,a4paper]{article}
\usepackage[a4paper,top=2.5cm,bottom=2.5cm,left=2.8cm,right=2.8cm,headheight=14pt]{geometry}
\usepackage{amsmath,amssymb}
\usepackage[T1]{fontenc}
\usepackage[utf8]{inputenc}
\usepackage{lmodern,microtype,fancyhdr,enumitem,hyperref,lastpage,parskip}

\hypersetup{
    colorlinks=true,
    urlcolor=black,
    linkcolor=black,
    pdftitle={BCF-215: Corporate Laws -- BHU B.Com. (Hons.) FMM Sem III 2013-14}
}

\pagestyle{fancy}
\fancyhf{}
\renewcommand{\headrulewidth}{0.4pt}
\renewcommand{\footrulewidth}{0.4pt}
\fancyhead[L]{\small   \textit{Banaras Hindu University}}
\fancyhead[R]{\small   \textit{BCF-215: Corporate Laws}}
\fancyfoot[C]{\small Page  \thepage\ of \pageref{LastPage}}


\newlist{parts}{enumerate}{2}
\setlist[parts,1]{label=(\alph*),leftmargin=*,itemsep=5pt,topsep=3pt}
\setlist[parts,2]{label=(\roman*),leftmargin=*,itemsep=3pt,topsep=2pt}

\newcommand{\pts}[1]{\hfill   \textit{[#1]}}


\begin{document}


\begin{center}
  {\large\bfseries BANARAS HINDU UNIVERSITY}\\[2pt]
  {
\normalsize B.Com. (Hons.) FMM --- Semester III Examination, 2013-14}\\[6pt]
  \rule{0.8\linewidth}{0.5pt}\\[6pt]
  {\large\bfseries Paper No. BCF-215: Corporate Laws}\\[4pt]
  {
\normalsize Time: 3 Hours\hfill Maximum Marks: 70}
\end{center}

\rule{\linewidth}{0.4pt}

\smallskip



\noindent   \textit{Instructions: Answer all questions. All questions carry equal marks (14 marks each).}

\bigskip




\noindent   \textbf{Question 1.} \\
``A company is an artificial person created by law having a separate entity with a perpetual succession and a common seal.'' Discuss this statement and explain the characteristics of a company.\pts{14}

\centerline{   \textbf{OR}}

Distinguish between the following:\pts{7+7}

\begin{parts}
  \item Private company and Public company
  \item Holding company and Subsidiary company
\end{parts}

\medskip\rule{\linewidth}{0.2pt}




\noindent   \textbf{Question 2.} \\
Explain the various stages involved with formation of a company under the Companies Act, 1956.\pts{14}

\centerline{   \textbf{OR}}

Discuss the functions and legal position of a promoter.\pts{14}

\medskip\rule{\linewidth}{0.2pt}




\noindent   \textbf{Question 3.} \\
What are the main clauses in Memorandum of Association of a company limited by shares? Describe the procedure of altering the object clause.\pts{14}

\centerline{   \textbf{OR}}

What is Articles of Association? Explain the importance and contents of Articles of Association.\pts{14}

\medskip\rule{\linewidth}{0.2pt}




\noindent   \textbf{Question 4.} \\
Discuss the qualifications and duties of director of a company.\pts{14}

\centerline{   \textbf{OR}}

Explain the legal position of directors of a company. Describe their powers and liabilities.\pts{14}

\medskip\rule{\linewidth}{0.2pt}




\noindent   \textbf{Question 5.} \\
How is the secretary in a company appointed? Discuss the duties of whole-time and part-time secretary.\pts{14}

\centerline{   \textbf{OR}}

Write short notes on any two of the following:\pts{7+7}

\begin{parts}
  \item Professionalisation of company secretaryship in India
  \item Issue and forfeiture of shares
  \item Dividend
\end{parts}

\vfill
\rule{\linewidth}{0.4pt}

\begin{center}\small
\href{https://bhu-exam.vercel.app/}{Click here to visit our website}\\[4pt]
\href{https://me-aryan.vercel.app/}{Click here to visit our contributor}\\[4pt]
\href{https://quashan-me.vercel.app/}{Click here to visit our contributor}
\end{center}

\end{document}
